\begin{Answer}{7.2}
(a) No.  The domain is $\BBZ$, which does not have a `starting point'.
(b) Yes.  The domain is $\BBZ^+$.
(c) Yes.  The domain is $\{2,3,4,\ldots\}$.
(d) Yes.  The domain is $\BBZ^+$.
(e) No.
The domain is $\reals$ which is not a subset of $\BBZ$.
Thus, not only is there no `starting point,'
there is no clear ordering of the real numbers from one to the next.
\end{Answer}
\begin{Answer}{7.4}
Modus ponens.
\end{Answer}
\begin{Answer}{7.5}
You can immediately conclude that $P(6)$ is true using
modus ponens.  If that was your answer, good.
But you can keep going.  Since $P(6)$ is true,
you can conclude that $P(7)$ is true (also by modus ponens).
But then you can conclude that $P(8)$ is true.
And so on.  The most complete answer you can give is that
$P(n)$ is true for all $n\geq 5$.
You {\em cannot} conclude that $P(n)$ is true for
all $n\geq 1$ because we don't know anything about
the truth values of $P(1)$, $P(2)$, $P(3)$, and $P(4)$.
\end{Answer}
\begin{Answer}{7.6}
Nothing. We can conclude that $P(n)$ is true for any $n\geq 17$,
but there is not enough information to say anything about values of $n$
less than 17.
\end{Answer}
\begin{Answer}{7.7}
There are various ways to say this, including what was said in the paragraph above.
Here is another way to say it:
\begin{quote}
If $P(a)$ is true, and for any value of $k\geq a$,
$P(k)$ true implies that $P(k+1)$ is true, then $P(n)$ is true for all $n\geq a$.
\end{quote}
\end{Answer}
\begin{Answer}{7.9}
If you answered {\bf yes} and you aren't lying, great!  If you
answered {\em no} or you answered {\bf yes} but you lied,
it is important that you think about it some more and/or
get some help.  If you want to succeed at writing induction
proofs, understanding this is an important step!
\end{Answer}
\begin{Answer}{7.12}
$\dfrac{1(1+1)}{2}$;
$P(1)$ is true;
$P(k)$ is true;
$\dsum_{i=1}^{k} i = \dfrac{k(k+1)}{2}$;
$P(k+1)$;
$\dsum_{i=1}^{k+1} i = \dfrac{(k+1)(k+2)}{2}$;
$\dsum_{i=1}^{k} i$;
$\dfrac{k(k+1)}{2}$;
$\dfrac{k}{2}+1$;
$\dfrac{(k+1)(k+2)}{2}$;
$P(k+1)$ is true;
$P(1)$ is true;
$k\geq 1$;
all $n\geq 1$;
induction {\em or} the principle of mathematical induction {\em or} PMI.
\end{Answer}
\begin{Answer}{7.13}
\begin{lenum}
\item $P(k)$ is the statement ``$\dsum_{i=1}^{k} i\cdot i! = (k+1)!-1$''
\item $P(k+1)$ is the statement ``$\dsum_{i=1}^{k+1} i\cdot i! = (k+2)!-1$''
\item $LHS(k)=\dsum_{i=1}^{k} i\cdot i!$
\item $RHS(k)=(k+1)!-1$
\item $LHS(k+1)=\dsum_{i=1}^{k+1} i\cdot i!$
\item $RHS(k+1)=(k+2)!-1$
\end{lenum}
\end{Answer}
\begin{Answer}{7.16}
Define $P(n)$ to be the statement ``$\dsum _{i=1}^n i^2 = \dfrac{n(n+1)(2n+1)}{6}$''.
We need to show that $P(n)$ is true for all $n\geq 1$.

\noindent
{\bf Base Case:}
Since $\dsum _{i=1}^1 i^2 = 1^2 = 1 =  \dfrac{1(2)(3)}{6}$,
$P(1)$ is true. (If your algebra is in a different order, like
$\dsum _{i=1}^1 i^2 = \dfrac{1(2)(3)}{6}= 1$, it is incorrect.  We only know that
$\dsum _{i=1}^1 i^2 = \dfrac{1(2)(3)}{6}$ because we first saw that $\dsum _{i=1}^1 i^2=1$,
and then were able to see that $1 =  \dfrac{1(2)(3)}{6}$.)

\noindent
{\bf Inductive Hypothesis:}
Let $k\geq 1$ and assume $P(k)$ is true.  That is, $\dsum _{i=1}^k i^2 = \dfrac{k(k+1)(2k+1)}{6}$.

\noindent
(As a side note, I know that what I need to prove next is

\noindent
$\displaystyle\dsum _{i=1}^{k+1} i^2 = \dfrac{(k+1)(k+2)(2(k+1)+1)}{6}=\dfrac{(k+1)(k+2)(2k+3)}{6}.$

\noindent
I am only writing this down now so that I know what my goal is.  I am not going to start working
both sides of this or otherwise manipulate it.  I can't because I don't know whether or not it is
true yet.)

\noindent
{\bf Inductive Step:}
Notice that
\begin{eqnarray*}
\dsum _{i=1}^{k+1} i^2&=&\dsum _{i=1}^k i^2+(k+1)^2\\
&=&\dfrac{k(k+1)(2k+1)}{6}+(k+1)^2\\
&=&(k+1)\left(\dfrac{k(2k+1)}{6}+(k+1)\right)\\
&=&(k+1)\left(\dfrac{k(2k+1)+6(k+1)}{6}\right)\\
&=&(k+1)\left(\dfrac{2k^2 + k + 6k+ 6}{6}\right)\\
&=&(k+1)\left(\dfrac{2k^2 + 7k+ 6}{6}\right)\\
&=&(k+1)\left(\dfrac{(2k+3)(k+2)}{6}\right)\\
&=&\dfrac{(k+1)(k+2)(2k+3)}{6}.
\end{eqnarray*}
Therefore $P(k+1)$ is true.

\noindent
{\bf Summary:} Since $P(1)$ is true and $P(k)\rightarrow P(k+1)$ is true when $k\geq 1$,
$P(n)$ is true for all $n\geq 1$ by induction.
\end{Answer}
\begin{Answer}{7.18}
For $k = 1$ we have $1\cdot 2 = 2 + (1-1)2^2$, and so the
statement is true for $n = 1$.
Let $k\geq 1$ and assume the statement is true for
$k$. That is, assume
$$1\cdot 2 + 2\cdot 2^2 + 3\cdot 2^3 + \cdots + k\cdot 2^k = 2 + (k - 1)2^{k + 1}.$$
We need to show that
$$1\cdot 2 + 2\cdot 2^2 + 3\cdot 2^3 + \cdots + (k+1)\cdot 2^{k+1} = 2 + k2^{k + 2}.$$
Using some algebra and the inductive hypothesis, we can see that
\begin{eqnarray*}
1\cdot 2 + 2\cdot 2^2 + 3\cdot 2^3 + \cdots + k\cdot 2^k + (k+1)2^{k+1}
&=& 2 + (k - 1)2^{k + 1} + (k+1)2^{k+1}\\
&=& 2 + (k - 1+k+1)2^{k + 1}\\
& =& 2 + 2k2^{k+1}\\
&=& 2 + k2^{k+2}.
\end{eqnarray*}
Thus, the result is true for $k+1$.  The result follows by induction.
\end{Answer}
\begin{Answer}{7.20}
This proof is very close to being correct, but is suffers from a few small but important errors:
\begin{itemize}
\item
For the sake of clarity, it might have been better to use $k$ throughout most of the proof
instead of $n$.  The exception is in the final sentence where $n$ is correct.
\item
The base case is just some algebra without context.  A few words are needed.  For instance,
`notice that when $n=1$,'.
\item
The base case is presented incorrectly.  Notice that the writer starts by writing down what
she wants to be true and then deduces that it is indeed correct by doing algebra on both sides
of the equation.  As we have already mentioned,
{\em you should {\bf never} start with what you want to prove and work both sides!}
It is not only sloppy, but it can lead to incorrect proofs.
Whenever I see students do this, I always tell them to use what I call the $U$ method.  What I mean
is rewrite your work by starting at the upper left, going down the left side, then doing up the
right side.  So the above should be rewritten as:
$$1\cdot 1! = 1 = 2!-1 = (1+1)!-1.$$
Notice that if the $U$ method does not work (because one
or more steps isn't correct), it is probably an indication
of an incorrect proof.  Consider what happens if you try it
on the proof in Exercise~\ref{minusOneIsNotOne}.  You would
write $-1=(-1)^2=1=1^2=1$.  Notice that the first equality
is incorrect.

The $U$ method can sometimes apply to inequalities as well.
\item
When the writer makes her assumption, she says `for $n\geq 1$'.  This is O.K., but there is some ambiguity
here.  Does she mean for {\em all} $n$, or for a particular value of $n$?  She must mean the latter since
the former is what she is trying to prove.  It would have been better for her to say `for some $n\geq 1$.'
\item
The algebra in the inductive step is perfect.  However, what does it mean?  She should include something like
`Notice that' before her algebra just to give it a little context.  It often doesn't take a lot of words,
but adding a few phrases here and there goes a long way to help a proof flow more clearly.
\item
She says `Therefore it is true for $n$'.  She must have meant $n+1$ since that is what she just proved.
\item
As with her assumption, her final statement could be clarified by saying `for all $n\geq 1$.'
\end{itemize}
Overall, the proof has almost all of the correct content.
Most of the problems have to do with presentation.
But as we have seen with
other types of proofs,
the details are really important to get right!
\end{Answer}
\begin{Answer}{7.21}
Given this proof, we know that $P(1)$ is true.
We also know that
$P(2)\rightarrow P(3)$, $P(3)\rightarrow P(4)$, etc,
are all true.
Unfortunately, the proof omits showing that $P(2)$ is true,
so modus ponens never applies.   In other words, knowing that
$P(2)\rightarrow P(3)$ is true does us no good unless we know
$P(2)$ is true, which we don't.  Because of this, we don't
know anything about the truth values of $P(3)$, $P(4)$, etc.
The proof either needs to show that $P(2)$ is true as part of
the base case, or the inductive step needs to start at $1$
instead of $2$.
\end{Answer}
\begin{Answer}{7.23}
Because our inductive hypothesis was that $P(k-1)$ is true instead of $P(k)$.
If we assumed that $k\geq 0$, then when $k=0$ it would mean we are assuming $P(-1)$ is true, and
we don't know whether or not it is since we never discussed $P(-1)$.
\end{Answer}
\begin{Answer}{7.27}
This contains a very subtle error.  Did you find it?
If not, go back and carefully re-read the proof and think
carefully--at least one thing said in the proof {\em must}
be incorrect.  What is it?

O.K., here it is:  The statement `goat 2 is in both collections'
is not always true.  If $n=1$, then the first collection
contains goats $1$ through $1$, and the second collection
contains goats $2$ through $2$.  In this case, there is no overlap
of goat $2$, so the proof falls apart.
\end{Answer}
\begin{Answer}{7.28}
\begin{pfevalenum}
\item
This solution is on the right track,  but it has several technical problems.
	\begin{itemize}
	\item
	The base case should
	be $k=0$, not $k=1$.
	\item
	The way the base case is worded could be improved.  For instance,
	what purpose does saying `$2=2$' serve?  Also, the separate sentence that just says `it is true' is
	a little vague and awkward.  I would reword this as:
		\begin{quote}
		The total number of palindromes of length $2\cdot 1$ is $2=2^1$, so the statement is true for $k=1$.
		\end{quote}
	Of course, the base case should {\em really} be $k=0$, but if it were $k=1$, that is how I would word it.
	\item
	The connection between palindromes of length $2k$ and $2(k+1)$ is not entirely clear and is incorrect
	as stated.  A palindrome of length $2(k+1)$ can be formed from a palindrome of length $2k$ by adding
	a $0$ to both the beginning and end or adding a $1$ to both the beginning and the end.
	This what was probably meant, but it is not what the proof actually says.
	
	But we need to say a little more about this.  Every palindrome of length
	$2(k+1)$ can be formed from exactly one palindrome of length $2k$ with this method.  But is this enough?
	Not quite.  We also need to know that every palindrome of length $2k$ can be extended to a palindrome of
	length $2(k+1)$, and it should be clear that this is the case.
	In summary, the inductive step needs to establish that there are twice as many binary palindromes of length $2(k+1)$ as there are of length $2k$.
	The argument has to convince the reader that there is a 2-to-1 correspondence between these sets of palindromes.  In other words, we did not omit or double-count any.
	\end{itemize}
\item
	The base case correct.  Unfortunately, that is about the only thing that is correct.
	\begin{itemize}
	\item
	The second sentence is wrong.
	We cannot say that `it is true for all $n$'--that is precisely what we are trying to prove.  We need
	to assume it is true for a {\em particular} $n$ and then prove it is true for $n+1$.
	\item
	The rest of the proof is one really long sentence that is difficult to follow.  It should be split
	into much shorter sentences, each of which provides one step of the proof.	
	\item
	The term `binary number' should be replaced with `binary palindrome' throughout.  It causes
	confusion, especially when the words `add' and `consecutive' are used.  These mean something
	very different if we have numbers in mind instead of strings.
	\item
	I don't think the phrase `each consecutive binary number' means what the writer thinks it means.
	The binary numbers $1001$ and $1010$ are consecutive (representing $9$ and $10$), but that is
	probably {\em not}	what the writer has in mind.
	\item
	The term `permutations' shows up for some reason.
	I think they might have mean `strings' or something else.
	\item
	Why bring up the $4$ possible ways to extend a binary
	string by adding to the beginning and end if only two of them are relevant?
	Why not just consider the ones of interest in the first place?
	\item
	In the context of a proof, the phrase 'you are adding' doesn't make sense.  Why am I adding something
	and what am I adding it to?  And do they mean addition (of the binary numbers) or appending (of strings)?
	\item
	They switch from $n$ to $k$ in the middle of the proof to provide further confusion.
	\end{itemize}
\item
This proof has most of the right ideas, but it does not put them together well.
The base case is correct.
It sounds like the writer understands what is going on with the inductive step,
but needs to communicate it more clearly.
More specifically, what does `assume $2k\rightarrow 2^k$ palindromes' mean?
I think I am supposed to read this as `assume that there are $2^k$ palindromes of length $2k$.'
\footnote{In general, avoid the use of mathematical symbols in constructing the grammar of an English sentence.
One of the most common abuses I see is the use of $\rightarrow$ in the middle of a sentence.}

The final sentence is also problematic.  The first phrase tries to connect to the previous sentence,
but the connection needs to be a little more clear.  The final phrase is not a complete thought.
In the first place, I know that  $2^k+2^k=2^{k+1}$ and this has nothing to do with the previous
phrases.  In other words, the `so' connecting the phrases doesn't make sense.  But more seriously,
why do I care that $2^k+2^k=2^{k+1}$?  What he meant was something like
`so there are $2^k+2^k=2^{k+1}$ palindromes of length $2k+2$'.
\end{pfevalenum}
\end{Answer}
\begin{Answer}{7.29}
The empty string is the only string of length $0$, and it is a palindrome.
Thus there is $1=2^0$ palindromes of length $0$.

Now assume there are $2^n$ binary palindromes of length $2n$.
For every palindrome of length $2n$, exactly two palindromes of length $2(n+1)$ can be constructed
by appending either a $0$ or a $1$ to both the beginning and the end.  Further, every palindrome
of length $2(n+1)$ can be constructed this way.  Thus, there are twice as many palindromes of
length $2(n+1)$ as there are of length $2n$.  By the inductive hypothesis, there are $2\cdot 2^n=2^{n+1}$
binary palindromes of length $2(n+1)$.

The result follows by PMI.
\end{Answer}
\begin{Answer}{7.32}
Yes.  It clearly calls itself in the else clause.
\end{Answer}
\begin{Answer}{7.35}
(a) The base cases are $n<\leq 0$.
(b) The inductive cases are $n>0$.
(c) Yes.  For any value $n>0$, the recursive call uses the value $n-1$, which is getting
closer to the base case of $0$.
\end{Answer}
\begin{Answer}{7.38}
Notice that if $n\leq 0$, \scode{countdown(0)} prints nothing, so it works in that case.
For $k\geq 0$, assume \scode{countdown(k)} works correctly.\footnote{We are letting $n=0$ be the base case.
You could also let $n=1$ be the base case, but then you would need to prove that \scode{countdown(1)} works.}
Then \scode{countdown(k+1)} will print `$k+1$'
and call \scode{countdown(k)}.  By the inductive hypothesis,
\scode{countdown(k)} will print `k\, k-1\, $\ldots$ 2\, 1', so \scode{countdown(k+1)} will print `k+1\, k\, k-1 $\ldots$ 2\, 1',
so it works properly.  By PMI,  \scode{countdown(n)} works for all $n\geq 0$.
\end{Answer}
\begin{Answer}{7.42}
It is pretty clear that the recursive algorithm
is much shorter and was a lot easier to write.
It is also a lot easier to
make a mistake implementing the iterative algorithm.
So far, it looks like the recursive algorithm is the clear winner.
However, in the next section we will show you why the recursive
algorithm we gave should never be implemented.  It turns out that is is
very inefficient.

The bottom line is that the iterative algorithm is better in this case.
Don't feel bad if you thought the recursive algorithm was better.
After the next section, you will be better prepared to compare
recursive and iterative algorithms in terms of efficiency.
\end{Answer}
\begin{Answer}{7.45}
\verb+PrintN+ will print from $1$ to $n$, and \verb+NPrint+ will print from $n$ to $1$.
If you go the answer wrong, go back and convince yourself that this is correct.
\end{Answer}
\begin{Answer}{7.48}
(a) $r_{n/2}$.
(b) $1$.
(c) $a_{n-1}+2\cdot a_{n-2}+3\cdot a_{n-3}+4\cdot a_{n-4}$.
(d) There are none.
\end{Answer}
\begin{Answer}{7.52}
It means to find a closed-form expression for it.  In other words, one that does not
define the sequence recursively.
\end{Answer}
\begin{Answer}{7.54}
When $n=1$, $T(1)=1=0+1=\log_2 1 + 1$.
Assume that $T(k)=\log_2 k + 1$ for all $1\leq k <n$
(we are using strong induction).
Then
\begin{eqnarray*}
T(n)&=&T(n/2)+1\\
&=&(\log_2 (n/2) + 1) + 1\\
&=&\log_2 n - \log_2 2 + 2\\
&=&\log_2 n - 1 + 2\\
&=&\log_2 n +1.
\end{eqnarray*}
So by PMI, $T(n)=\log_2 n + 1$ for all $n\geq 1$.
\end{Answer}
\begin{Answer}{7.56}
We begin by computing a few values to see if we can find a pattern.
$A(2)=A(1)+2 = 2+2= 4$, $A(3)=A(2)+2=4+2=6$, $A(4)=8$, $A(5)=10$, etc.
It seems pretty obvious that $A(n)=2n$.  It holds for $n=1$, so we have our base case.
Assume $A(n)=2n$.  Then $A(n+1)=A(n)+2 = 2n+2=2(n+1)$, so it holds for $n+1$.
By PMI, $A(n)=2n$ for all $n\geq 1$.
\end{Answer}
\begin{Answer}{7.59}
It contains 3 very different looking recursive terms so it is very unlikely we will
be able to find any sort of meaningful pattern by iteration.
\end{Answer}
\begin{Answer}{7.61}
\begin{eqnarray*}
H(n)&=&2 H(n-1) + 1\\
&=&2(2 H(n-2) + 1) + 1\\
&&= 2^2 H(n-2) + 2 + 1\\
&=&2^2(2 H(n-3) + 1) + 2 + 1\\
&&= 2^3 H(n-3) + 2^2 + 2 + 1\\
&\vdots&\\
&=&2^{n-1}H(1) + 2^{n-2} + 2^{n-3}+\cdots + 2 + 1\\
&=&2^{n-1} + 2^{n-2} + 2^{n-3}+\cdots + 2 + 1\\
&=&2^n-1
\end{eqnarray*}
Thus, $H(n)=2^n-1$. Luckily, this matches our answer from Example~\ref{exa:Hanoi1}.
\end{Answer}
\begin{Answer}{7.63}
Iterating a few steps, we discover:
\begin{eqnarray*}
T(n)&=&T(n/2)+1\\
&=&T(n/4)+1+1\\
&=&T(n/2^2)+2\, \, \, \, \mbox{(I think I see a pattern!)}\\
&=&T(n/2^3)+1+2\\
&=&T(n/2^3)+3\, \, \, \, \mbox{(I {\em do} see a pattern!)}\\
&\vdots&\\
&=&T(n/2^k)+k\\
\end{eqnarray*}
We need to find $k$ such that $n/2^k=1$.
We already saw in Example~\ref{exa:longiterone} that $k=\log_2 n$ is the solution.
Therefore, we have
\begin{eqnarray*}
T(n)&=&T(n/2^k)+k\\
&=&T(n/2^{\log_2 n})+\log_2 n\\
&=&T(1)+\log_2 n\\
&=&1+\log_2 n\\
\end{eqnarray*}
Therefore,
$T(n)=1+\log_2 n$.
\end{Answer}
\begin{Answer}{7.64}
The final answer is $T(n)=2^{n+1}-n-2$ or $T(n)=4\cdot2^{n-1}-n-2$.
It is important that you can work this out yourself, so try your best to get this
answer without looking further.  But if you get stuck or have a different answer,
you can refer to the following skeleton of steps--we have omitted many of the steps
because we want you to work them out.  It does provide a few reference points
along the way, however.
\begin{eqnarray*}
T(n)&=&2T(n-1)+n\\
&=&2(2T(n-2)+(n-1))+n\, \, \, \mbox{(having $n$ instead of $(n-1)$ is a common error)}\\
&=&2^2T(n-2)+3n-2\, \, \, \mbox{(it is unclear yet if I should have $3n-2$ or some other form)}\\
&\vdots& \, \, \, \, \, \, \mbox{(many skipped steps)}\\
&=&2^kT(n-k)+(2^k-1)n-\dsum_{i=1}^{k-1}i2^i\, \, \, \, \mbox{(the all-important pattern revealed)}\\
&\vdots& \, \, \, \, \, \, \mbox{(plug in appropriate value of $k$ and simplify)}\\
&=&2^{n+1}-n-2.
\end{eqnarray*}
\end{Answer}
\begin{Answer}{7.69}
Here, $a=2, b=2$, and $d=0$. ($d=0$ since $1=1\cdot n^0$.  In general, $c=c\cdot n^0$, so when
$f(n)$ is a constant, $d=0$.)
Since $a>2^0$, we have $T(n)=\Theta(n^{\log_aa})=\Theta(n^1)=\Theta(n)$
by the third case of the Master Theorem.
\end{Answer}
\begin{Answer}{7.71}
We have $a=7$, $b=2$, and $d=2$.  Since $7>2^2$, the third case of the
Master Theorem applies so $T(n)=\Theta(n^{\log_{2} 7})$, which is about $\Theta(n^{2.8})$.
\end{Answer}
\begin{Answer}{7.72}
Because it isn't true.  Although the growth rate of $n^{\log_{2} 7}$ and $n^{2.8}$ are
close, they are not exactly the same, so $\Theta(n^{\log_{2} 7})\neq \Theta(n^{2.8})$.
We {\em could} say that $T(n)=\Theta(n^{\log_{2} 7})=O(n^{2.81})$, but then we have lost
the `tightness' of the bound.  And I want to be able to say ``Yo dawg, that bound is really {\em tight}!''
\end{Answer}
\begin{Answer}{7.73}
Here we have $a=1$, $b=2$, and $d=0$.  Since $1=2^0$,
the second case of the Master Theorem tells is that $T(n)=\Theta(n^0\log n)=\Theta(\log n)$.
Since we have already seen several times that $T(n)=\log_2 n +1$, we can notice that
this answer is consistent with those.  It's a good thing.
\end{Answer}
\begin{Answer}{7.79}
By raising the subscripts in the homogeneous equation we
obtain the characteristic equation $x^n = 9x^{n - 1}$ or $x = 9$.
A solution to the homogeneous equation will be of the form $x_n =
A(9)^n$. Now $f(n) = -56n + 63$ is a polynomial of degree 1 and so
we assume that the solution will have the form $x_n = A9^n + Bn + C$.
Now $x_0 = 2,
x_1 = 9(2) - 56 + 63 = 25, x_2 = 9(25) - 56(2) + 63 = 176$. We
thus solve the system
$$2 = A + C,$$
$$25 = 9A + B + C,$$
$$176 = 81A + 2B + C.$$
We find $A = 2, B = 7, C = 0$, so the solution is $x_n = 2(9^n) + 7n.$
\end{Answer}
\begin{Answer}{7.83}
The characteristic equation is $x^2 - 4x + 4 = (x - 2)^2
= 0$. There is a multiple root and so we must test a solution of
the form $x_n = A2^n + Bn2^n.$ The initial conditions give
$$1 = A,$$
$$4 = 2A + 2B.$$
This solves to $A = 1, B = 1.$ The solution is thus $x_n = 2^n + n2^n$.
\end{Answer}
\begin{Answer}{7.87}
We have $a=2$, $b=2$, and $d=1$. Since $2=2^1$, we have that
$T(n)=\Theta(n\log n)$ by the second case of the Master Theorem.
\end{Answer}
\begin{Answer}{7.88}
(a) $C_1$ and $C_2$ are of lower order than $n$.
Thus, we can do this according to part (c) of Theorem~\ref{thm:asyprops}.
(b) We know that the $\Theta(n)$ represents some function $f(n)$.  By definition of $\Theta$,
there are constants, $c_1$ and $c_2$ such that $c_1n\leq f(n)\leq c_2n$ for all $n\geq n_0$
for some constant $n_0$.  So we essentially replaced $f(n)$ with $c_2n$ (since we are
looking for a worst-case).
Note that it doesn't matter what $c_2$ is.
We know there is some constant that works, so we just call it $c$.
\end{Answer}
\begin{Answer}{7.90}
$T(n)=2 T(n-1)+T(n-5)+T(\sqrt{n})+1$ if $n>5$,
$T(1)=T(2)=T(3)=T(4)=T(5)=1$.
You can also have $+c$ instead of $+1$ in the recursive definition.
\end{Answer}
\begin{Answer}{7.91}
Beyond the recursive calls, \verb+StoogeSort+ does only a constant amount of work--we'll call it $1$.
Then it makes three calls with sub-arrays of size $(2/3)n$.  Therefore,
$T(n)=3T((2/3)n)+1$ or $T(n)=3T(2n/3)+1$, with base case $T(1)=1$.
\end{Answer}
\begin{Answer}{7.92}
We can use the Master Theorem for this one.
$a=3$, $b=3/2$ and $d=0$.
(Notice that $b\neq 2/3$!  If you made this mistake, make sure you understand why it is incorrect.)
Since $3>1=(3/2)^0$, the third case of the Master Theorem tells us that
$T(n)=\Theta\left(n^{\log_{3/2}(3)}\right)$.  Although this looks weird, we can have a rational number as
the base of a logarithm (in fact, the base of $ln(n)$ is $e$, an irrational number).
It might be helpful to compute the $\log$ since it isn't clear how good or bad this complexity is.
Notice that $\log_{3/2}(3)\approx 2.71$, so $T(n)$ is approximately $\Theta(n^{2.71})$, but
$T(n)\neq \Theta(n^{2.71})$, so resist the urge to place an equals sign between these.
\end{Answer}
\begin{Answer}{7.93}
The complexity of \verb+MergeSort+ is $\Theta(n\log n)$ and the complexity of
\verb+StoogeSort+ is $\Theta\left(n^{\log_\frac{3}{2}(3)}\right)$
which is $\Omega(n^{2.7})$.
Clearly $\Theta\left(n^{\log_\frac{3}{2}(3)}\right)$ grows faster than $\Theta(n\log n)$,
so \verb+MergeSort+ is faster.
{\em Remember: Faster growth rate means slower algorithm!}
\end{Answer}
\begin{Answer}{7.94}
Let $T(n)$ be the complexity of this algorithm.
From the description, it seems pretty clear that $T(n)=5T(n/3)+cn$.
Using the Master Theorem with $a=5$, $b=3$, and $d=1$, we see that $5>3^1$,
so the third case applies and $T(n)=\Theta\left(n^{\log_3(5)}\right)$,
which is approximately $\Theta(n^{1.46})$.
\end{Answer}
